\textbf{1. Kiến thức cần nhớ}
\begin{itemize}
\item Với $A$, $B$ là các biểu thức tùy ý:
\begin{itemize}
\item Lập phương của một tổng: $\boxed{(A+B)^3=A^3+3A^2B+3AB^2+B^3}$.
\item Lập phương của một hiệu: $\boxed{(A-B)^3=A^3-3A^2B+3AB^2-B^3}$.
\end{itemize}
\item Nhận xét: khai triển có 4 hạng tử, hệ số $1;\,3;\,3;\,1$; lũy thừa của $A$ giảm dần, của $B$ tăng dần. Với lập phương của một hiệu, dấu đan xen $+,-,+,-$.
\item $(A-B)^3=-(B-A)^3$.
\end{itemize}

\textbf{2. Phương pháp giải các dạng}

%% DANG 8D2?7-1
\textbf{Dạng 1. Khai triển lập phương của một tổng, một hiệu}
\begin{enumerate}
\item Xác định $A$, $B$; tính sẵn $A^2$, $A^3$, $B^2$, $B^3$.
\item Thay vào hằng đẳng thức và thu gọn: $(2x+1)^3=8x^3+12x^2+6x+1$.
\end{enumerate}
Lưu ý: $(A+B)^3\ne A^3+B^3$; lập phương của số âm là số âm.

%% DANG 8D2?7-2
\textbf{Dạng 2. Viết biểu thức dưới dạng lập phương}
\begin{enumerate}
\item Sắp xếp đa thức theo lũy thừa giảm dần; tìm hạng tử đầu $A^3$ và hạng tử cuối $\pm B^3$.
\item Kiểm tra hai hạng tử giữa có bằng $\pm 3A^2B$ và $3AB^2$ không.
\item Đúng thì viết $(A+B)^3$ (dấu toàn $+$) hoặc $(A-B)^3$ (dấu đan xen).
\end{enumerate}

%% DANG 8D2?7-3
\textbf{Dạng 3. Tính nhanh, rút gọn biểu thức}
\begin{enumerate}
\item Tính nhanh: đưa biểu thức về $(A\pm B)^3$ với $A\pm B$ là số tròn, ví dụ $x^3+6x^2+12x+8=(x+2)^3$.
\item Rút gọn: khai triển các lập phương, bỏ ngoặc (đổi dấu khi trước ngoặc có dấu $-$), thu gọn hạng tử đồng dạng.
\end{enumerate}
Lưu ý: tính giá trị biểu thức thì rút gọn trước rồi mới thay số.
